\documentclass[../main.tex]{subfiles}
\begin{document}

\FloatBarrier
\section{Summary}
\label{sec:summary}

This thesis studied four-class ordinal student-engagement recognition from per-frame OpenFace facial-behaviour descriptors and clip-level I3D motion features.
It set out to improve recognition over standard monolithic models and, equally, to characterise the problem honestly --- in particular its severe, heterogeneous class imbalance and its behaviour under distribution shift.

Methodologically, the thesis built a complete and reproducible pipeline: feature extraction and leakage-free preprocessing; a \textbf{semantic-group hybrid temporal architecture} that decomposes the 709-dimensional OpenFace descriptor into five behavioural groups (gaze, eye landmarks, face landmarks, head pose, action units), encodes each group with its own temporal encoder, and optionally fuses a global I3D stream while supervising every stream with an auxiliary head; five monolithic baselines; three loss functions; and a 3$\times$3 cross-dataset evaluation protocol.
The central idea was validated by training and evaluating all $3^5=243$ per-group encoder assignments, with and without I3D, and by a \textbf{self-collected, test-only private set} of 366 hand-labelled clips that serves as the decisive deployment probe.

The evidence supports five conclusions.
\begin{enumerate}
  \item \textbf{The I3D-fused hybrid is the best in-domain model and, as a family, beats the strongest monolithic baseline.}
  The I3D-fused family is centred above the baseline QWK of 0.537 (median 0.553, with 83\% of its 243 configurations clearing the baseline), and the best configuration reaches \textbf{QWK 0.605} (accuracy 0.784) with a mixed TCN/Transformer/LSTM assignment (\texttt{TCN\_T\_TCN\_LSTM\_T}).
  The gain over the best baseline is genuine but \textbf{modest} ($+0.068$ QWK); the value lies in the framework and the systematic evidence, not in a leap in state of the art.
  \item \textbf{Most groups are encoder-agnostic; only head pose shows a clear preference, favouring a TCN} (a 0.022 QWK spread across the three encoders, against $\le0.006$ for the other four groups).
  This is an interpretable design rule reflecting the temporally local nature of head dynamics.
  \item \textbf{Engagement classifiers do not generalise across datasets.}
  Off-diagonal QWK collapses to near zero (CMOSE$\rightarrow$DAiSEE 0.02) even when accuracy looks acceptable, so ordinal agreement metrics --- not accuracy --- must be reported on imbalanced engagement data.
  \item \textbf{On the self-collected, hand-labelled private set --- the most realistic test --- the two contributions compound.}
  Because the private set is test-only, the only lever is the training source; training on the \textbf{combined} corpus with the \textbf{semantic-group hybrid} gives the best result (QWK 0.379), beating the best baseline by $+0.094$ (against $+0.068$ in-domain) and topping every single-source hybrid.
  This is the practical payoff of the thesis.
  \item \textbf{The rare Highly\nobreakdash-Disengage class remains the principal open problem.}
  It is around 2--3\% of the data, is traded away rather than solved by the hybrid in-domain, and collapses to zero recall under distribution shift.
\end{enumerate}

A note on scope: the architecture ablation is anchored on CMOSE because DAiSEE's in-domain ordinal signal is too weak (QWK 0.166) for architecture differences to be distinguishable from label noise; CMOSE is therefore used for development and the private set for the decisive validation.

\FloatBarrier
\section{Synthesis}
\label{sec:synthesis}

Read together, the two results chapters tell a single story whose moral is not the one a benchmark paper usually reports.
The thesis has a \emph{constructive} side --- the semantic-group hybrid and the corpus-pooling recipe --- and a \emph{critical} side --- an honest account of imbalance, non-transfer, and the metrics that expose them --- and the most important finding lies in how the two interact.

\textbf{The architecture earns its value off-domain, not on it.}
In-domain the semantic-group decomposition is a modest, interpretable improvement: the I3D-fused family sits above the baseline but the best configuration clears it by only $+0.068$ QWK, and most of the 243-cell design space is encoder-agnostic.
Taken alone, that is a small result.
What makes it worth the construction is that the same decomposition pays off \emph{more} exactly where it is hardest to: on the unseen private set the margin over the baseline widens to $+0.094$, and it compounds with pooling rather than competing with it.
The contribution is therefore best understood not as a new in-domain state of the art but as a representation that degrades more gracefully under distribution shift --- which, for a problem whose only real test is deployment, is the property that matters.

\textbf{The flat design space is itself the message.}
That four of five behavioural groups are indifferent to their encoder, while only head pose prefers a TCN, bounds the contribution honestly: the gain comes from \emph{decomposing and fusing} the descriptor, not from a clever per-group encoder search.
This is a more transferable design rule than a single tuned configuration, and it is the reason the family-level effect, rather than any one cell, is reported as the result.

\textbf{The negative results are contributions, not caveats.}
The near-total collapse of cross-corpus transfer, the way raw accuracy disguises it, and the persistent failure on the rarest disengaged class are not incidental limitations of these particular models; they are properties of the problem that the field's in-domain, accuracy-led habit leaves unmeasured.
Establishing them quantitatively --- and showing that QWK and macro-accuracy make them visible --- reframes engagement recognition from ``beat the benchmark'' to ``report ordinal agreement and instrument rare-class failure,'' which is the discipline the rest of this chapter recommends.

\FloatBarrier
\section{Revisiting the Research Questions}
\label{sec:rq-revisit}

The four research questions posed in Section~\ref{sec:objectives} can now be answered directly.

\textbf{RQ1 (decomposition).}
\emph{Yes, decomposition helps.}
Encoding the five behavioural groups separately and fusing them outperforms encoding all 709 features with a single monolithic encoder: as a family the hybrid is centred above the best baseline, and the effect is robust rather than a single lucky configuration, with most of the I3D-fused configurations clearing the baseline.

\textbf{RQ2 (encoder assignment).}
\emph{Only head pose matters.}
Four of the five groups are essentially indifferent to whether their encoder is a TCN, a Transformer, or an LSTM (spread $\le0.005$ QWK); head pose alone shows a clear and interpretable preference for a TCN.
The practical implication is that the design space is mostly flat and that a TCN is a safe default for every group, with head pose the one place the choice is worth making deliberately.

\textbf{RQ3 (motion fusion).}
\emph{Yes, the I3D stream helps and stabilises the gain.}
Adding the clip-level I3D embedding shifts the whole family of configurations upward and produces both the best in-domain model and the best private-set model, even though the I3D feature is only a globally pooled clip descriptor.

\textbf{RQ4 (generalisation and losses).}
\emph{Generalisation is poor and metric choice is decisive.}
Models do not transfer off-the-shelf across corpora; off-diagonal ordinal agreement collapses while accuracy stays deceptively high.
Among losses there is no universal winner --- cross-entropy maximises accuracy while weighted and ordinal losses recover minority recall --- so the loss is a Pareto choice, and QWK together with macro-accuracy must be the primary yardstick.
Pooling source corpora is the single most effective lever for an unknown target.

\FloatBarrier
\section{Limitations}
\label{sec:limitations}

The findings should be read against several limitations, each of which also motivates the future work below.
\begin{itemize}
  \item \textbf{Single-annotator private labels.}
  The private set was labelled by one annotator (the author) under the public-corpus rubric.
  Although the model suggestions shown during labelling were non-binding, the absence of a second annotator means no inter-rater reliability can be reported for this set, and the 366-clip size limits the statistical power of the private-set conclusions --- especially for the ten Highly\nobreakdash-Disengage clips, where a single misjudged clip moves the recall substantially.
  \item \textbf{Clip-level I3D features.}
  CMOSE and DAiSEE ship a single pooled 1024-dimensional I3D vector per clip, so the I3D ``stream'' is an MLP over that global motion/appearance summary rather than a temporal encoder; a per-window I3D sequence might contribute more and would let an I3D temporal encoder model genuine dynamics.
  \item \textbf{Incremental in-domain gain and a flat design space.}
  Most of the 243-configuration grid is encoder-agnostic, and the best in-domain improvement over the baseline is modest ($+0.068$ QWK).
  The decomposition is a sound, interpretable framework rather than a large accuracy gain.
  \item \textbf{Weak DAiSEE signal.}
  The very low in-domain QWK on DAiSEE prevented using it for architecture development and means the cross-dataset conclusions rest mainly on CMOSE and the private set.
  \item \textbf{Fixed training recipe.}
  A single seed, learning rate, and encoder size were used throughout; the absence of multi-seed repetition means the reported margins are point estimates rather than confidence intervals, and small differences between neighbouring configurations should not be over-interpreted.
\end{itemize}

\FloatBarrier
\section{Practical Recommendations}
\label{sec:recommendations}

For practitioners building engagement-monitoring systems, the study yields three concrete recommendations that are independent of the specific architecture.
First, \textbf{report and optimise ordinal agreement}: use QWK and macro-accuracy as the headline metrics and treat raw accuracy as secondary, because on these imbalanced ordinal labels accuracy can look healthy while the model has learned nothing useful about disengagement.
Second, \textbf{pool heterogeneous source corpora} when the deployment distribution is unknown; pooling was the most effective simple lever for the unseen private set and requires no change to the model.
Third, \textbf{expect and instrument rare-class failure}: the most decision-relevant class, Highly\nobreakdash-Disengage, is also the rarest and the least reliably detected, so a deployed system should treat low-confidence disengagement predictions cautiously and, where possible, defer to a human.

\FloatBarrier
\section{Suggestion for Future Works}
\label{sec:future}

\begin{itemize}
  \item \textbf{Targeting the rarest class.}
  Combine the architecture with imbalance-specific remedies --- ordinal losses tuned for Highly\nobreakdash-Disengage, minority oversampling such as SMOTE~\cite{smote}, or cost-sensitive decision thresholds --- to recover HD recall without trading away Highly\nobreakdash-Engage.
  Because HD is so rare, even modest gains here would change the practical value of the system more than further QWK gains on the frequent classes.
  \item \textbf{Domain adaptation.}
  The cross-dataset gap invites explicit unsupervised domain adaptation or domain-invariant representation learning, rather than the naive corpus pooling that already helps; a small amount of labelled target data could be used for calibration before deployment.
  \item \textbf{Learned, not enumerated, group encoders.}
  Replace the discrete 243-configuration search with a differentiable architecture search or a gating mechanism that selects per-group temporal operators end-to-end, removing the need to enumerate the design space and allowing the model to discover group--encoder affinities automatically.
  \item \textbf{Stronger and richer motion features.}
  Re-extract per-window I3D (or a modern video backbone) so the motion stream carries temporal dynamics, and fuse audio/speech features, which CMOSE also provides, as additional semantic streams within the same decompose-then-fuse framework.
  \item \textbf{Frequency-domain analysis of engagement (exploratory).}
  A natural complementary direction --- not implemented in this thesis and left strictly as future work --- is to analyse the \emph{temporal frequency} of facial signals, under the hypothesis that engaged behaviour is low-frequency and stable while disengaged behaviour is higher-frequency and fidgety, for example via short-time-Fourier-transform features or convolutions over the frequency axis.
  \item \textbf{Stronger evaluation.}
  Multi-annotator labelling of a larger private set, with reported inter-rater agreement, and multi-seed training with confidence intervals would tighten every margin reported here and turn point estimates into statistically grounded claims.
\end{itemize}

\FloatBarrier
\section{Closing Remarks}
\label{sec:closing}

Automatic engagement recognition is attractive precisely because online learning makes the webcam the only window onto the learner, but the same setting makes the problem hard: the labels are ordinal, severely imbalanced, subjective, and corpus-specific.
This thesis showed that respecting the structure of the facial descriptor --- dividing it into behavioural groups and encoding each on its own terms --- yields a sound and interpretable model that, fused with a motion stream and trained on pooled corpora, generalises best to the most realistic, unseen data.
Just as importantly, it showed that the field's habitual reliance on accuracy hides a near-total failure to transfer across datasets, and that ordinal agreement metrics expose it.
Recognising the rarest, most decision-relevant disengaged behaviour under distribution shift remains open, and is, in the author's view, the most worthwhile next step.

\end{document}
